This work uses two openly licensed, non-credentialed datasets, each serving
a distinct role in the pipeline: one as the corpus the retriever searches
over, the other as the held-out benchmark the downstream task is graded on.
Table~\ref{tab:datasets} summarizes both; full preprocessing detail is in
the project's \texttt{docs/dataset.md}.

\begin{table}[t]
\centering
\caption{Datasets used in this work.}
\label{tab:datasets}
\begin{tabular}{@{}lll@{}}
\toprule
Role & Dataset & License \\
\midrule
Retrieval corpus & ROCOv2~\cite{ruckert2024rocov2} & CC BY-NC-SA 4.0 \\
Evaluation benchmark & VQA-RAD~\cite{lau2018vqarad} & CC0 (public domain) \\
\bottomrule
\end{tabular}
\end{table}

\subsection{ROCOv2: Retrieval Corpus}

ROCOv2~\cite{ruckert2024rocov2} pairs approximately 90{,}000 radiological
images with captions extracted from the PMC Open Access subset, spanning
multiple modalities (CT, MRI, X-ray, and ultrasound). Each caption is a
clinically written figure caption from published biomedical literature,
not a crowd-sourced description. This corpus plays the role of the
non-parametric evidence store in the retrieval-augmented condition
(Section~\ref{sec:method}): at inference time, the query image is embedded
and used to retrieve the most visually similar corpus entries, whose
captions are then supplied to the model as evidence.

One caveat carried through to the retrieval design: ROCOv2 captions are
figure captions, not full clinical reports, and occasionally reference
comparative or positional context (e.g., a prior figure panel) that is not
recoverable from the isolated retrieved image. Section~\ref{sec:discussion}
returns to whether this measurably affects retrieved-evidence quality.

\subsection{VQA-RAD: Evaluation Benchmark}

VQA-RAD~\cite{lau2018vqarad} contains 3{,}515 question-answer pairs over 315
radiology images, with questions written directly by clinicians rather than
crowdworkers, covering both closed-ended (yes/no-style) and open-ended
free-text questions across CT, MRI, and X-ray studies. Using a benchmark
disjoint from the retrieval corpus is deliberate: it evaluates the
downstream VQA task independently of the retriever's own source corpus,
rather than measuring the model's ability to recover an answer key it was
given evidence from.

\subsection{Splitting and Preprocessing}

VQA-RAD is evaluated on its own official held-out test split (451
examples), rather than a further internal re-split of the training
portion, so that results are directly comparable to other published
VQA-RAD baselines that use the same standard split. The official
training portion (1{,}793 examples) is further divided 90/10 into
train/validation internally, stratified by answer type (closed
vs.\ open) where that label is available. The Hugging Face mirror of
VQA-RAD used here (\texttt{flaviagiammarino/vqa-rad}) does not expose
an explicit closed/open answer-type field; this label is instead
recovered from the answer text itself, treating a normalized answer of
exactly ``yes'' or ``no'' as closed-ended and every other answer as
open-ended. This is an approximation of the original VQA-RAD
annotation -- a small number of closed-ended, non-yes/no questions
(e.g.\ ``which side is the lesion on?'') will be mislabeled
open-ended -- and is noted as a limitation in
Section~\ref{sec:limitations}. ROCOv2 is not split, since it serves
only as the full retrieval corpus, not a held-out evaluation set.
Images are resized to $224\times224$ and normalized with standard
ImageNet statistics, noted here as a documented default rather than a
finding specific to medical images, which have different intensity
distributions than natural images; captions are truncated to 128
whitespace-delimited tokens before being written to the retrieval
manifest.

\subsection{Why Not MIMIC-CXR}

MIMIC-CXR is the chest radiograph dataset most commonly used in published
medical VLM work. It was not used here because it requires a credentialed
PhysioNet access process (CITI training and institutional affiliation
review) with a latency of one to several weeks, which would have delayed
every downstream milestone of a timeline-constrained project. ROCOv2 and
VQA-RAD require no credentialing and no protected health information
handling, allowing the full pipeline to be built and validated end to end
first. Extending the retrieval corpus to include MIMIC-CXR, once
credentialed access is obtained, is identified as future work
(Section~\ref{sec:limitations}).
